\documentclass[11pt]{article}
\usepackage[utf8]{inputenc}
\usepackage[T1]{fontenc}
\usepackage[margin=2.3cm]{geometry}
\usepackage{amsmath,amssymb}
\usepackage{booktabs}
\usepackage{graphicx}
\usepackage{adjustbox}
\usepackage{xcolor}
\usepackage{float}
\usepackage{caption}
\usepackage[hidelinks]{hyperref}
\usepackage{textcomp}

\setlength{\parskip}{4pt}
\renewcommand{\figurename}{Figura}
\renewcommand{\tablename}{Tabla}
\renewcommand{\refname}{Referencias}

% ----------------------------------------------------------------------
% Las cifras, tablas y figuras las genera el notebook en tablas/ y figuras/.
% Si falta alguna (porque el notebook no se ha corrido), sale en rojo.
% ----------------------------------------------------------------------
\InputIfFileExists{tablas/cifras.tex}{}{}
\newcommand{\pendiente}{\textcolor{red}{\textbf{[correr notebook]}}}
\newcommand{\cf}[1]{\ifcsname #1\endcsname\csname #1\endcsname\else\pendiente\fi}
\newcommand{\tabla}[1]{\IfFileExists{tablas/#1.tex}{\begin{adjustbox}{max width=\textwidth}\input{tablas/#1.tex}\end{adjustbox}}{\pendiente}}
\newcommand{\figura}[2][0.95]{\IfFileExists{figuras/#2.png}{\includegraphics[width=#1\textwidth]{figuras/#2.png}}{\pendiente}}
% \revisar marca frases que dependen de cómo salgan sus datos: léanlas contra el notebook y quiten el color.
\newcommand{\revisar}[1]{\textcolor{orange!80!black}{#1}}

\title{Tarea 4: Introducción a Series de Tiempo\\[2pt]\large Business Intelligence · Universidad del Istmo de Guatemala}
\author{José Sierra}
\date{29 de octubre de 2026}

\begin{document}
\maketitle

\noindent Todo el código está en \texttt{Tarea4\_Sierra.ipynb}; la limpieza, en \texttt{limpiar.py}, y las
funciones de ACF, Ljung--Box y AR(1), implementadas con numpy, en \texttt{ts\_utils.py}. El notebook corre de principio a
fin desde \texttt{Online Retail.xlsx} y escribe las figuras en \texttt{figuras/}, las tablas en \texttt{tablas/} y
las cifras que se citan en el texto en \texttt{tablas/cifras.tex}, así que este documento se actualiza solo al
volver a compilarlo.

%======================================================================
\section{Parte 0: Procesos simulados}
%======================================================================

\subsection*{0.1 Derivaciones a mano}

Sea $Y_t = c + \phi Y_{t-1} + \varepsilon_t$ con $\varepsilon_t \overset{iid}{\sim} N(0,\sigma^2)$, $|\phi|<1$ y el
proceso estacionario, de modo que $E[Y_t]=E[Y_{t-1}]=\mu$ y $\operatorname{Var}(Y_t)=\operatorname{Var}(Y_{t-1})=\gamma_0$.

\paragraph{Media.} Tomando esperanza a ambos lados y usando $E[\varepsilon_t]=0$:
\[
\mu = c + \phi\mu \quad\Longrightarrow\quad \mu(1-\phi)=c \quad\Longrightarrow\quad E[Y_t]=\frac{c}{1-\phi}.
\]

\paragraph{Varianza.} Como $\varepsilon_t$ es independiente del pasado, en particular de $Y_{t-1}$,
la covarianza entre ambos es cero y
\[
\gamma_0 = \operatorname{Var}(c+\phi Y_{t-1}+\varepsilon_t) = \phi^2\gamma_0 + \sigma^2
\quad\Longrightarrow\quad \operatorname{Var}(Y_t)=\frac{\sigma^2}{1-\phi^2}.
\]

\paragraph{Autocorrelación.} Restando la media, $Y_t-\mu = \phi(Y_{t-1}-\mu)+\varepsilon_t$. Multiplicando por
$Y_{t-k}-\mu$ con $k\ge 1$ y tomando esperanza, el término $E[\varepsilon_t(Y_{t-k}-\mu)]$ es cero porque
$\varepsilon_t$ es independiente de $Y_{t-k}$:
\[
\gamma_k = \phi\,\gamma_{k-1} \;\Longrightarrow\; \gamma_k = \phi^k\gamma_0 \;\Longrightarrow\; \rho_k=\frac{\gamma_k}{\gamma_0}=\phi^k.
\]

\paragraph{Caminata aleatoria.} Con $\phi=1$, $c=0$ y $Y_0=0$, iterando, $Y_t=\sum_{j=1}^t\varepsilon_j$. Como los
$\varepsilon_j$ son independientes,
\[
\operatorname{Var}(Y_t)=\sum_{j=1}^t\operatorname{Var}(\varepsilon_j)=t\sigma^2.
\]
La varianza crece con $t$, así que no es constante en el tiempo y se viola la definición de estacionariedad débil.
Además $\operatorname{Cov}(Y_t,Y_{t-k})=(t-k)\sigma^2$ depende de $t$ y no solo de $k$. Por eso las fórmulas de
arriba no aplican: con $\phi=1$, $\sigma^2/(1-\phi^2)$ ``explota''.

\subsection*{0.2 Simulación}

Cada trayectoria tiene $n=300$ pasos después de descartar 200 de calentamiento, con $c=0$ y $\sigma=1$. La
figura~\ref{fig:real} muestra una realización de cada $\phi$: con $\phi=0$ la serie cruza la media constantemente;
con $\phi=0.9$ hace rachas largas, y con $\phi=1$ deambula sin volver a ningún nivel.

\begin{figure}[H]\centering\figura[0.85]{p0_realizaciones}
\caption{Una realización de cada AR(1).}\label{fig:real}\end{figure}

Para $\phi=1$ el calentamiento no borra el valor inicial (la caminata nunca olvida), pero eso no afecta a $r_k$:
el valor de arranque suma la misma constante a todas las observaciones, y $r_k$ resta la media muestral antes
de multiplicar, así que la constante se cancela.

La tabla~\ref{tab:02} resume 5\,000 réplicas por escenario, con una semilla distinta para cada uno
(documentadas en el notebook). EE es el error estándar del promedio: la desviación estándar entre réplicas
dividida entre $\sqrt{5000}$ (para la tasa de rechazo, $\sqrt{p(1-p)/5000}$).

\begin{table}[H]\centering\small\tabla{p0_tabla02}
\caption{AR(1) con $n=300$ y 5\,000 réplicas por $\phi$.}\label{tab:02}\end{table}

\subsection*{0.3 Interpretación}

\paragraph{$\phi=0$.} La tasa de rechazo es \cf{rechazoPhiCero}, a menos de un error estándar del 5\,\% que
esperábamos: bajo ruido blanco $r_1\approx N(0,1/n)$, así que $|r_1|>1.96/\sqrt n$ ocurre el 5\,\% de las veces por
construcción. El promedio de $r_1$ es ligeramente negativo ($\approx -1/n$), un sesgo conocido de $r_k$ por usar la
media muestral.

\paragraph{Caminata aleatoria.} La teoría ``estacionaria'' diría $\rho_5=1^5=1$, pero el promedio de $r_5$ es
\cf{rCincoPhiUno}. La fórmula $\rho_k=\phi^k$ supone estacionariedad, que aquí no existe: $\bar y$ no es la media
de nada, y en una muestra finita $r_k$ siempre queda por debajo de 1 y decae despacio. Esa caída lenta es justamente
la pista visual de no estacionariedad.

\paragraph{Sesgo del OLS.} El estimado de $\phi$ es sistemáticamente menor que el verdadero, y el sesgo crece con
$\phi$: \cf{sesgoPhiCero}, \cf{sesgoPhiMedio}, \cf{sesgoPhiNueve} y \cf{sesgoPhiUno} para $\phi=0,\ 0.5,\ 0.9,\ 1$.
Los errores estándar son de $10^{-3}$ o menos, así que el sesgo no es ruido de simulación. Coincide con la
aproximación de Kendall, $-(1+3\phi)/n$ ($-0.0033$, $-0.0083$, $-0.0123$), y para $\phi=1$ con el sesgo de
Dickey--Fuller con intercepto, cerca de $-5.3/n\approx-0.018$. El estimador es consistente (el sesgo es de orden
$1/n$ y desaparece cuando $n$ crece), pero en muestras finitas subestima la memoria, más cuanto más persistente es
el proceso.

\subsection*{0.4 Las dos razones para reprobar}

\begin{table}[H]\centering\small\tabla{p0_tabla04}
\caption{Procesos sin memoria con media cambiante; 5\,000 réplicas.}\label{tab:04}\end{table}

\begin{figure}[H]\centering\figura[0.95]{p0_dos_razones}
\caption{Arriba: una realización y su media $m_t$. Abajo: ACF promedio de 5\,000 réplicas con bandas $\pm1.96/\sqrt n$.}\label{fig:dos}\end{figure}

\paragraph{La tendencia reprueba sin tener memoria.} Rechaza la prueba de retardo 1 en el \cf{rechazoUnoTendencia}
de las réplicas, con $\bar r_1=\cf{rUnoTendencia}$, aunque los $\varepsilon_t$ son independientes. Los días de la
parte alta de la tendencia quedan todos sobre $\bar y$ y sus vecinos también, así que el producto
$(y_t-\bar y)(y_{t-1}-\bar y)$ es positivo casi siempre. Su ACF (figura~\ref{fig:dos}) baja muy despacio: es la
razón (2).

\paragraph{El patrón semanal pasa aunque tiene estructura.} Rechaza solo en el \cf{rechazoUnoSemanal} (poco más del
5\,\% nominal) porque al rezago 1 casi no hay correlación ($\bar r_1=\cf{rUnoSemanal}$). La estructura vive en el
rezago 7: $\bar r_7=\cf{rSieteSemanal}$. La prueba mira el rezago equivocado.

\paragraph{Conexión con RANDU.} RANDU pasaba las pruebas de pares consecutivos y aun así todas sus ternas
$(u_t,u_{t+1},u_{t+2})$ caen en 15 planos. En los dos casos hay una estructura determinista fuerte en una
dimensión o un rezago que la prueba no mira. Pasar una prueba solo dice que no hay estructura \emph{donde la
prueba busca}.

\subsection*{0.5 La prueba que sí ve}

\begin{table}[H]\centering\small\tabla{p0_tabla05}
\caption{Tasa de rechazo al 5\,\% de Ljung--Box con $m=10$ y con $m=\log n\approx\cf{mLog}$.}\label{tab:05}\end{table}

Con ruido blanco, Ljung--Box con $m=10$ rechaza el \cf{lbDiezRuido}: mantiene su tamaño. Ahora el proceso semanal
se rechaza en el \cf{lbDiezSemanal} de las réplicas (contra \cf{rechazoUnoSemanal} con la prueba de retardo 1),
y la tendencia en el \cf{lbDiezTendencia}. Ljung--Box suma $r_k^2$ de los rezagos 1 a 10, así que incluye el
rezago 7, donde está la estructura semanal, y una sola $r_7$ grande basta para que $Q$ sea grande. Además, al
juntar 10 rezagos en un estadístico evita el problema de las pruebas múltiples.

El enunciado pide comparar con la regla de Tsay, $m\approx\log n=\cf{mLog}$: con $m=6$ el proceso semanal solo se
rechaza en el \cf{lbLogSemanal}, porque el rezago 7 queda fuera. Es el argumento de FPP2 para usar $m=2s$ con
estacionalidad: $m$ debe cubrir al menos el periodo estacional.

%======================================================================
\section{Parte 1: Limpieza}
%======================================================================

Leímos \texttt{InvoiceNo}, \texttt{StockCode} y \texttt{CustomerID} como texto. Leídos como números,
\texttt{InvoiceNo} mezcla \texttt{536365} con \texttt{C536379} y la columna queda con tipos mezclados;
\texttt{StockCode} mezcla \texttt{85123} con \texttt{85123A}, y \texttt{CustomerID}, al tener vacíos, se convierte en
\texttt{float} (\texttt{17850.0}), lo que rompe los cruces y los conteos por cliente.

\subsection*{1.1 Tamaño}

\begin{table}[H]\centering\tabla{p1_tamano}\caption{La base cruda.}\end{table}

El \cf{fracSinCliente} de las filas no tiene \texttt{CustomerID}.

\subsection*{1.2 Casos y decisiones}

\begin{description}
\item[(a) Duplicados exactos (\cf{nDuplicados} filas): se van.] Dos filas idénticas en factura, producto, cantidad,
precio y hora son un doble registro del sistema; dejarlas contaría dos veces la misma venta. Se conserva la primera.

\item[(b) Facturas que no empiezan con dígito: C se queda, A se va.] Hay dos prefijos (tabla~\ref{tab:pref}).
\texttt{C} son cancelaciones: se quedan y restan el día en que ocurren, como pide la definición de venta neta.
\texttt{A} son ajustes contables \emph{Adjust bad debt}: no son ventas a un cliente, así que se van.

\item[(c) Códigos que no son productos: se van.] La tabla~\ref{tab:cod} muestra los códigos que no empiezan con cinco
dígitos. Por su descripción, \texttt{POST}, \texttt{DOT} y \texttt{C2} son cobros de envío; \texttt{AMAZONFEE},
\texttt{BANK CHARGES} y \texttt{CRUK} son comisiones y cargos; \texttt{M} (Manual), \texttt{D} (Discount),
\texttt{S} (Samples) y \texttt{B} son ajustes, y \texttt{gift\_...} son vales de regalo. Ninguno es venta de
mercancía y se van. En cambio \texttt{PADS} y los \texttt{DCGS...} sí son productos y se quedan, igual que los
códigos con letra al final (\texttt{85123A}), que son variantes de un producto.

\item[(d) Precio cero o negativo (\cf{nPrecioCero} y \cf{nPrecioNeg} filas): se van.] Una línea con precio cero no
aporta venta. El 98.4\,\% de ellas no tiene cliente y sus descripciones son vacías o del tipo \emph{check},
\emph{damages}, \emph{found}, \emph{adjustment}: son movimientos de inventario. Los dos precios negativos son los
ajustes de deuda incobrable (facturas A) de (b).

\item[(e) Cantidades negativas fuera de cancelaciones (\cf{nQNeg} filas): se van.] Las descripciones son del tipo
\emph{damaged}, \emph{lost}, \emph{?}, \emph{check}: son correcciones de inventario, no devoluciones de clientes.
El \cf{pctQNegPrecioCero} de ellas tiene precio cero y ninguna tiene cliente, así que la regla (d) ya las elimina
todas y la regla 6 queda en cero filas.

\item[(f) Sin \texttt{CustomerID}: para la serie se quedan; para la Parte 6 se van.] Representan el
\cf{pctMontoSinCliente} de la venta bruta y son ventas reales, así que sacarlas sesgaría la serie hacia abajo.
Para RFM sí se excluyen, porque necesitamos saber quién compró.
\end{description}

Además agregamos una regla que no estaba en la lista: \textbf{pedidos anulados con un ajuste Manual}. Si un pedido
grande se cancela con una línea \texttt{M} del mismo cliente y el mismo monto, la regla (c) borra la cancelación y
deja el pedido vivo, con lo que la limpieza \emph{crearía} una venta que nunca existió. Por eso, antes de (c),
quitamos esas líneas de venta (ver 1.4).

\begin{table}[H]\centering\small\tabla{p1_prefijos}\caption{Facturas por prefijo.}\label{tab:pref}\end{table}
\begin{table}[H]\centering\scriptsize\tabla{p1_codigos}\caption{Códigos que no empiezan con cinco dígitos (los 25 más frecuentes).}\label{tab:cod}\end{table}

\subsection*{1.3 Bitácora de limpieza}

Las reglas se aplican en el orden de la tabla~\ref{tab:bit}. El orden importa: la regla 6 no elimina ninguna fila porque la regla 5 ya quitó todas las cantidades negativas
(todas tenían precio cero), y si se aplicara antes se llevaría 1\,336 filas que hoy cuenta la regla 5. Del mismo modo,
la regla 3 tiene que ir antes de la 4: después ya no existiría la línea Manual que delata al pedido anulado. El monto con signo dice cuánto
cambia la venta neta; el monto absoluto, cuánta actividad se tocó.

\begin{table}[H]\centering\tabla{p1_bitacora}\caption{Bitácora de limpieza.}\label{tab:bit}\end{table}

Después de limpiar quedan \cf{nFilasLimpias} filas con una venta neta de \cf{ventaNetaLimpia}.

\subsection*{1.4 Las seis líneas más grandes}

\begin{table}[H]\centering\tabla{p1_top6}\caption{Las seis líneas con mayor $|Quantity\times UnitPrice|$ antes de limpiar
y su contrapartida.}\label{tab:top6}\end{table}

Las seis líneas son tres pares pedido--anulación. El cliente 16446 pidió 80\,995 unidades de \emph{Paper Craft,
Little Birdie} (factura 581483, £168\,469.60) el 9 de diciembre a las 9:15 y la canceló 12 minutos después (C581484).
El cliente 12346 pidió 74\,215 \emph{Medium Ceramic Top Storage Jar} (541431, £77\,183.60) el 18 de enero y lo
canceló 16 minutos después (C541433). Son los ``pedidos que no deberían existir'': cantidades absurdas para una
tienda cuyo pedido típico es de decenas de unidades. Nuestra limpieza no los elimina (pedido y cancelación son
líneas válidas), pero netear por día sí los resuelve: la venta y la cancelación caen el mismo día y se anulan. El
neto del 18 de enero queda en £17\,067, un día normal, y el 9 de diciembre además se descarta en 2.4.

El tercer par es distinto: el cliente 15098 pidió 60 \emph{Picnic Basket Wicker 60 Pieces} a £649.50 (556444,
£38\,970) y la anulación no se registró con el producto, sino con una línea \texttt{M} (Manual) de £38\,970
(C556445). Aquí netear por día no basta si la regla 4 borra el Manual: el 10 de junio quedaría con £38\,970 de
venta fantasma. Por eso la regla 3 quita ese pedido; con ella el neto del 10 de junio es £21\,937, un día normal.

\subsection*{1.5 Función \texttt{limpiar}}

Toda la limpieza es la función \texttt{limpiar(df)} de \texttt{limpiar.py}. Devuelve el DataFrame limpio con
columnas \texttt{Venta}, \texttt{Fecha} y \texttt{EsCancelacion}; con \texttt{devolver\_bitacora=True} devuelve
también la bitácora, y con \texttt{solo\_clientes=True} quita las filas sin cliente para la Parte 6.

%======================================================================
\section{Parte 2: El calendario}
%======================================================================

\paragraph{2.1} Entre el primer y el último día de datos, incluyendo ambos, hay \cf{diasCalendario} días calendario,
y hubo al menos una transacción en \cf{diasConTx}.

\paragraph{2.2} Los \cf{diasCerrados} días sin transacciones se clasifican en la tabla~\ref{tab:cerr}. La tienda
\textbf{no abre los sábados}: los \cf{nSabados} sábados del rango están vacíos. Los otros 16 días son tres grupos. El
\textbf{cierre de Navidad}, del viernes 24 de diciembre de 2010 al lunes 3 de enero de 2011 (reabre el 4 de enero),
que incluye los \emph{bank holidays} sustitutos de Christmas Day, Boxing Day y Año Nuevo. La \textbf{Semana Santa}:
Good Friday, Easter Monday y también el domingo de Pascua, a pesar de que la tienda sí abre los domingos. Y los
\textbf{bank holidays de lunes} (2 y 30 de mayo, 29 de agosto) más el viernes 29 de abril, feriado por la boda
real del príncipe William: ese es el que no se repite cada año.

Los móviles en el sentido de FPP2 son los tres de Semana Santa, porque dependen de la fecha de Pascua: en 2011
cayeron del 22 al 25 de abril, y en 2010 o 2012 caen en otras semanas, así que un modelo de datos diarios no puede
aprenderlos con ``el mismo día del año pasado''. Los de mayo y agosto también cambian de fecha, pero siempre son
lunes, y la boda real es un evento único que ningún modelo puede anticipar con estos datos.

\begin{table}[H]\centering\small\tabla{p2_cerrados}\caption{Días sin transacciones que no son sábado.}\label{tab:cerr}\end{table}

\paragraph{2.3 ¿Cero o faltante?} Depende del uso. Debe ser \textbf{cero} cuando se quiere el total vendido en un
periodo (venta mensual, presupuesto, flujo de caja) o una media móvil de calendario como la de 3.1: ese día
realmente no entró dinero. Debe ser \textbf{faltante} cuando se modela la venta de los días en que la tienda opera,
por ejemplo para estimar el efecto del día de la semana o evaluar un pronóstico: un sábado cerrado no es un día de
ventas bajas, y tratarlo como cero arrastraría el índice y castigaría a cualquier modelo que no ``prediga'' el
cierre.

\paragraph{2.4 El último día.} La última factura del 9 de diciembre es a las \cf{horaUltima}, mientras que la mediana
de la hora de la última factura en los días abiertos es \cf{horaCierreMediana}: el día está incompleto (los datos
se cortaron a mediodía), aunque su venta neta de la mañana ya sea parecida a la de un viernes completo. Lo quitamos,
porque compararlo con un día completo haría parecer que las ventas se desplomaron.

%======================================================================
\section{Parte 3: Tendencia, estacionalidad y autocorrelación}
%======================================================================

\subsection*{3.1 Tendencia}

\begin{figure}[H]\centering\figura{p3_tendencia}\caption{Venta neta diaria y media móvil centrada de 7 días.}\label{fig:tend}\end{figure}

Usamos 7 días porque es el periodo de la estacionalidad: cada ventana contiene exactamente un día de cada día de
la semana, así que el patrón semanal se promedia y desaparece. Con 5 días las ventanas incluirían unas veces el
sábado cerrado y otras no, y la media oscilaría con el patrón semanal; con 30 se borraría el patrón, pero también
los cambios de la tendencia en el mes (por ejemplo la subida de noviembre). Es centrada para que el valor de cada
día sea un promedio de tres días antes y tres después, sin desfasar la tendencia; una media hacia atrás llegaría
tarde a cada cambio.

Diciembre de 2010 arranca alto (la media móvil está cerca de £40\,000 diarios), cae a cero en el cierre de
Navidad y de enero a agosto se mueve en un rango plano de unos £15\,000 a £30\,000 diarios, con subidas cortas
como la de mayo. Desde septiembre sube con fuerza y en noviembre pasa de £50\,000: es la temporada en que los
mayoristas de artículos de regalo se abastecen para Navidad. En total, noviembre vendió £1.43 millones, casi el
triple que febrero (£0.50 millones), así que la serie no tiene media constante.

\subsection*{3.2 Estacionalidad}

\begin{table}[H]\centering\tabla{p3_indice}\caption{Índice estacional: promedio de $y_t/MA_t$ en los días abiertos del entrenamiento.}\label{tab:ind}\end{table}
\begin{figure}[H]\centering\figura[0.65]{p3_indice}\caption{Índice estacional por día de la semana.}\end{figure}

El día con mayor índice es el \cf{diaMaxIndice} y el de menor, el \cf{diaMinIndice}. Jueves (\cf{indiceJue}) y martes (\cf{indiceMar})
venden más de 1.5 veces la media móvil; el domingo vende apenas \cf{indiceDom}, menos de la mitad que un día entre
semana. Es un día corto y con pocos pedidos mayoristas. El promedio de los seis índices es
\cf{promIndice} y no 1 porque la media móvil incluye el sábado cerrado con valor cero: el promedio de 7 días es
aproximadamente $6/7$ del nivel de un día abierto, así que los cocientes de los días abiertos promedian cerca de
$7/6\approx1.167$. Los feriados empujan en el mismo sentido: una ventana con un feriado tiene dos ceros. Sin las
ventanas con más de un día cerrado el promedio es \cf{promIndiceLimpio}.

\subsection*{3.3 Autocorrelación}

\begin{figure}[H]\centering\figura[0.8]{p3_acf}\caption{ACF de la venta en los \cf{nTres} días abiertos del entrenamiento (una fila por día).}\label{fig:acf3}\end{figure}

\paragraph{(a)} Con $n=\cf{nTres}$ las bandas son $\pm\cf{bandaTres}$, y se salen \cf{fueraTres} de los 21 rezagos.
$r_1=\cf{rUnoTres}$, así que la serie no pasa la prueba de retardo 1 de la Tarea 2. El pico está en el rezago
\cf{picoTres} ($r=\cf{rPicoTres}$). Como en esta serie los sábados no son filas, 6 filas son 7 días calendario:
el pico es la estacionalidad semanal. Ljung--Box da $Q(6)=\cf{QseisTres}$ ($p=\cf{pseisTres}$) y
$Q(12)=\cf{QdoceTres}$ ($p=\cf{pdoceTres}$): se rechaza ruido blanco con cualquier nivel razonable.

\paragraph{(b)} Por la Parte 0 sabemos que una ACF así tiene dos posibles fuentes. Los picos en los rezagos 6, 12 y 18
son la estacionalidad semanal, como el $r_7$ del proceso semanal de la Parte 0.4. Y que casi todos los demás rezagos
sean positivos, en especial los primeros, es la huella de una media no constante: la tendencia de la
figura~\ref{fig:tend} hace que los días de septiembre a noviembre queden todos por encima de $\bar y$, igual que en
el proceso con tendencia de la figura~\ref{fig:dos}. Ninguna de las dos es memoria. Lo que quede después de
quitarlas sí será memoria, y eso es lo que mide la Parte 4.

%======================================================================
\section{Parte 4: Del AR(1) al SARIMA}
%======================================================================

\subsection*{4.1 Quitar el día de la semana}

\begin{figure}[H]\centering\figura[0.8]{p4_acf_dow}\caption{ACF de los residuos de la regresión contra los seis indicadores de día.}\end{figure}

El pico en el rezago 6 desaparece: antes era el mayor ($r_6=\cf{rPicoTres}$) y ahora es uno más, porque las
indicadoras capturan el nivel promedio de cada día. Pero siguen fuera de las bandas \cf{fueraCuatroUno} rezagos, y la
forma no es la de un AR(1). Con $r_1=\cf{rUnoCuatroUno}$, un AR(1) daría $r_2\approx0.08$ y $r_3\approx0.02$; en
cambio, los rezagos 1 a 6 se quedan todos entre 0.21 y 0.30, y solo después caen. Es la meseta de la ACF del proceso
con tendencia de la Parte 0.4: semanas enteras de días por encima (o por debajo) del promedio, porque el nivel de la
serie se mueve. Al modelo le falta la tendencia (el nivel que se
mueve en el año); mientras esté en los residuos, cualquier prueba va a reprobar por la razón (2).

\subsection*{4.2 AR(1) sobre los residuos}

\begin{figure}[H]\centering\figura[0.8]{p4_acf_ar1}\caption{ACF de los residuos $u_t$ del AR(1).}\end{figure}

El AR(1) da $\hat\phi=\cf{phiAR}$. Sus residuos $u_t$ siguen con rezagos fuera de las bandas
(\cf{rezagosFueraAR}), y Ljung--Box con un grado de libertad menos da $Q(6)=\cf{QseisAR}$ ($p=\cf{pseisAR}$, 5 gl) y
$Q(12)=\cf{QdoceAR}$ ($p=\cf{pdoceAR}$, 11 gl). No alcanzó: los residuos siguen sin ser ruido blanco. El AR(1) solo
explica la correlación de un día con el anterior, pero lo que vimos en 4.1 es un nivel que persiste durante
semanas, y eso deja correlación en rezagos más largos, incluido el estacional 6: el efecto de cada día no es un
monto fijo, sino que crece con el nivel de la serie.

\subsection*{4.3 SARIMA}

\begin{table}[H]\centering\scriptsize\tabla{p4_sarima}\caption{Convergencia y parámetros de S1 y S2.}\end{table}

\paragraph{¿Por qué $s=6$ y no 7?} Porque ajustamos sobre las filas de días abiertos: sin sábados, una semana son 6
filas y el mismo día de la semana anterior está 6 filas atrás. La suposición se rompe en las semanas con feriados o
con el cierre de Navidad: ahí la semana tiene 5 filas o menos, y el rezago 6 ya no cae en el mismo día de la semana.

\paragraph{¿Por qué el logaritmo?} La amplitud del patrón semanal crece con el nivel de la serie (las diferencias
entre jueves y domingo son mayores en noviembre que en febrero): la estacionalidad es multiplicativa, como en la
descomposición de 3.1--3.2. El logaritmo la vuelve aditiva y estabiliza la varianza, que es cuando FPP2 recomienda
transformar. \revisar{Revisamos la convergencia en \texttt{mle\_retvals}; para el modelo que no convergió subimos
\texttt{maxiter} a 500 y la tabla compara los parámetros antes y después.}

\subsection*{4.4 Diagnóstico}

\begin{table}[H]\centering\tabla{p4_diagnostico}\caption{Ljung--Box sobre los residuos (con $m-2$ grados de libertad).}\end{table}
\begin{figure}[H]\centering\figura[0.95]{p4_acf_sarima}\caption{ACF de los residuos de S1 y S2.}\end{figure}

¿Residuos de ruido blanco al 5\,\%? S1: \cf{ruidoBlancoS1}; S2: \cf{ruidoBlancoS2}. \revisar{S2 deja residuos mucho
más cercanos a ruido blanco, porque sus dos diferencias quitan el nivel que se mueve y el patrón semanal; S1
modela el logaritmo alrededor de una media fija y la tendencia de otoño le queda en los residuos.}

\subsection*{4.5 Horizonte largo}

\begin{figure}[H]\centering\figura{p4_horizonte}\caption{Pronósticos de S1 y S2 a 120 días abiertos desde el 30 de septiembre.}\end{figure}

S1 es estacionario: como un AR(1) con $|\phi|<1$, olvida sus choques a ritmo $\phi^h$, así que su pronóstico regresa a
la media incondicional (con el patrón semanal amortiguándose) y a horizonte largo es una línea plana en el nivel
promedio de todo el entrenamiento. S2 tiene raíces unitarias: como la caminata aleatoria, no regresa a ningún
nivel. Su pronóstico repite indefinidamente el último nivel y el último patrón semanal estimados. Ninguno de los
dos ``sabe'' que viene la temporada navideña, pero S2 al menos parte del nivel reciente.

%======================================================================
\section{Parte 5: Pronóstico y fuga}
%======================================================================

\subsection*{5.1 Modelos}

B1 es la media histórica; B2, el último valor del mismo día de la semana; M1, OLS con tendencia lineal e
indicadoras; M2, OLS con indicadoras, $y_{t-7}$ e $y_{t-14}$. Para M2, si $t-7$ cae en un día cerrado tomamos el mismo
día de la semana abierto más reciente; $y_{t-14}$ es el mismo día de la semana abierto inmediatamente anterior a ese.
Comprobamos el caso del enunciado: para el 9 de mayo, $y_{t-7}$ es el 18 de abril.

\paragraph{¿Por qué $y_{t-7}$ sí e $y_{t-1}$ no?} El origen es un domingo y los objetivos están a lo sumo 7 días
después, así que para cualquier objetivo $t$ el día $t-7$ es el origen o anterior: ya está observado. En cambio, para
pronosticar el miércoles, $y_{t-1}$ sería el martes, que todavía no ha pasado. Usarlo sería usar el futuro.

\paragraph{¿Cómo usa S1 el rezago 1?} Pronostica recursivamente: el pronóstico del lunes se usa como si fuera el
valor observado para pronosticar el martes, y así sucesivamente. Por eso el horizonte de cada objetivo es su
posición entre los días abiertos y el error crece con el horizonte. En cada reajuste de S1 y S2 usamos el
\texttt{maxiter} de 4.3 y verificamos \texttt{mle\_retvals['converged']}; \revisar{\cf{noConvergio} ajustes no
convergieron}.

\subsection*{5.2 Tres formas de estimar el error}

\begin{table}[H]\centering\tabla{p5_errores}\caption{MAE y RMSE (en £). (a) validación cruzada aleatoria de 5 folds,
media $\pm$ desviación estándar en 10 semillas; (b) origen móvil semanal, mayo--septiembre; (c) prueba, origen
móvil semanal.}\label{tab:err}\end{table}

\paragraph{¿Por qué no (a) para B2, S1 y S2?} Porque los tres dependen del orden de la serie: B2 necesita ``el último''
valor, y S1 y S2 son recursivos sobre filas consecutivas. Si se sacan días al azar, la historia de entrenamiento
queda con huecos en cualquier lugar y el ``pasado'' de un día de prueba incluye días que vienen después de él. El
modelo no se puede ajustar ni pronosticar de forma coherente.

\paragraph{¿Cuál se acerca más a (c)?} Para M2 la más cercana es \cf{masCercaM2} (MAE de M2: (a)
\cf{maeAM2}, (b) \cf{maeBM2}, (c) \cf{maeCM2}). Pero ninguna se acerca de verdad: para B1, M1 y M2 el
error de prueba es entre 55\,\% y 120\,\% mayor que cualquiera de las dos estimaciones dentro del entrenamiento,
porque octubre--diciembre está en otro nivel de ventas.

\subsection*{5.3 ¿Quién gana en prueba?}

\begin{figure}[H]\centering\figura{p5_prueba}\caption{Pronósticos semanales contra la venta real, 3 de octubre -- 8 de diciembre de 2011.}\end{figure}

Gana \cf{ganadorPrueba} (MAE de \cf{maeCB1}, \cf{maeCB2}, \cf{maeCM1}, \cf{maeCM2}, \cf{maeCS1} y \cf{maeCS2} para
B1, B2, M1, M2, S1 y S2).

Lo que separa a los modelos en prueba es si se adaptan al nivel reciente. La venta de octubre--diciembre gira
alrededor de £50\,000 diarios. B1 se queda en la media histórica de los días abiertos
(entre £27\,700 y £31\,400 según el origen) y es por mucho el peor. M1 y M2 siguen bien el patrón semanal (la figura muestra que suben y bajan
con la serie), pero quedan sistemáticamente por debajo de ella: M1 extrapola una tendencia lineal estimada sobre
meses casi planos, y en M2 los coeficientes de $y_{t-7}$ e $y_{t-14}$ estimados en el entrenamiento suman apenas 0.23
(0.17 y 0.06), así que el pronóstico le hace poco caso a la semana anterior y regresa hacia el promedio de cada día. B2 copia la semana anterior completa: arrastra el nivel nuevo y por eso
tiene el menor error entre los modelos sin SARIMA, a pesar de que copia también los picos aislados (el día de
£110\,600 del 14 de noviembre lo repite el 21, cuando la venta real fue de £47\,000).

\revisar{En cuanto a S1 y S2: S1 es estacionario y, como vimos en 4.5, regresa a la media del año, así que debería
comportarse como B1 y quedar entre los peores. S2, con sus dos diferencias, se ancla al nivel reciente como B2,
pero promedia las semanas anteriores en vez de copiar solo la última, y modela el patrón semanal en escala
multiplicativa; si en 4.4 deja residuos de ruido blanco, es el candidato natural a ganar.}

\subsection*{5.4 La trampa}

\begin{table}[H]\centering\tabla{p5_trampa}\caption{M2 con y sin la media móvil centrada de 7 días como predictor.}\end{table}

Con la media móvil, M2 ``mejora'' en (a) y en (b) (MAE (b): \cf{maeBM2} sin ella, \cf{maeBTrampa} con ella). Es una
ilusión: la media móvil centrada del día $t$ promedia $y_{t-3},\dots,y_{t+3}$, así que contiene al propio $y_t$ y a
tres días futuros. El predictor le pasa al modelo parte de la respuesta.

El origen móvil no lo detecta si la media móvil se calcula una sola vez sobre toda la serie, porque el origen
móvil solo controla qué \emph{filas} entran al ajuste, no qué información hay \emph{dentro} de cada columna. La
columna ya fue construida con datos posteriores al origen, y la fuga entra por ahí.

Y (c) no se puede calcular sin hacer trampa: el domingo de origen, para pronosticar el jueves, necesitaríamos la
media móvil del jueves, que requiere las ventas del lunes al domingo siguiente, que todavía no ocurren. Además, la
media móvil de los últimos días de diciembre requiere datos que ni siquiera están en la base.

\subsection*{5.5 Por qué todas las estimaciones se quedan cortas}

Hay dos causas. La primera viene de \textbf{cómo se partieron los datos}. La validación cruzada aleatoria pone en
entrenamiento días vecinos a los de prueba (el martes y el jueves de la misma semana), así que el modelo ``ve'' el
nivel de esa semana; además, los rezagos de M2 pueden venir de días que en la realidad todavía no ocurrieron. Esa
causa \emph{se arregla} con una mejor validación: el origen móvil (b) respeta el tiempo y por eso es más honesto
que (a).

La segunda es que \textbf{la serie cambió}. Octubre--diciembre es la temporada alta de un mayorista de artículos de
regalo: el nivel de ventas es más alto y más volátil que en cualquier mes del entrenamiento. Ninguna validación
dentro del entrenamiento puede medir un error en condiciones que nunca observó; eso \emph{no se arregla} con una
mejor partición, solo con datos de otras temporadas altas.

Con un solo año no se puede estimar la estacionalidad anual. Para estimar un patrón con periodo de un año
necesitamos verlo repetirse: con un año hay una sola observación de cada mes, y es imposible distinguir ``noviembre
siempre es alto'' de ``la tienda está creciendo'' o de ``este noviembre fue raro''.

%======================================================================
\section{Parte 6: Detrás de la serie hay clientes: RFM}
%======================================================================

\subsection*{6.1 R, F y M al 1 de octubre}

Con los datos limpios que tienen \texttt{CustomerID} y solo las transacciones anteriores al 1 de octubre de 2011,
hay \cf{nClientesRFM} clientes con al menos una factura de venta. La recencia son los días entre su última factura
de venta y el 1 de octubre; la frecuencia, el número de facturas de venta distintas, y el monto, la venta neta
(ventas menos cancelaciones). \cf{nSoloCanc} clientes solo tienen cancelaciones antes del corte (de compras
anteriores al inicio de los datos) y no entran, porque no tienen recencia.

\cf{nMcero} clientes tienen $M\le0$: cancelaron tanto o más de lo que compraron. Entre ellos está el cliente 12346 del pedido de 74\,215 frascos
anulado en 1.4, que queda con $M=0$. Los conservamos: el enunciado pide incluirlos en los puntajes, y
quedan naturalmente en el quintil más bajo de $M$. No serían candidatos a una campaña de ``mejores clientes''.

\begin{table}[H]\centering\small\tabla{p6_describe}\caption{Distribución de R (días), F (facturas) y M (£).}\end{table}

No usamos nada posterior al 1 de octubre porque en 6.3 evaluamos si el RFM anticipa quién compra después de esa
fecha. Si el RFM usara compras de octubre, un cliente que compró en octubre tendría recencia baja \emph{porque}
compró: estaríamos prediciendo el resultado con el resultado (la misma fuga de 5.4).

\subsection*{6.2 Puntajes por quintiles}

F es un conteo con muchísimos empates: el \cf{pctFuno} de los clientes tiene exactamente una factura. Los puntos de
corte de los quintiles coinciden (el primer y el segundo quintil son ambos $F=1$), y \texttt{pd.qcut} falla por
bordes repetidos. Desempatamos al azar como pide el enunciado: clientes ordenados por \texttt{CustomerID},
barajados con \texttt{default\_rng(2026)} y \texttt{qcut} sobre \texttt{rank(method='first')}; lo mismo para R
(con signo invertido, para que 5 sea el más reciente) y M.

Eso implica que dos clientes con la misma F pueden quedar con puntajes distintos solo por el orden del barajado. La
diferencia de puntaje entre ellos no dice nada de su comportamiento; es ruido que agregamos para que los quintiles
tengan el mismo tamaño. Cambiar la semilla cambiaría a quién le toca cada puntaje.

\subsection*{6.3 ¿Qué dimensión separa mejor?}

\begin{table}[H]\centering\small\tabla{p6_tasa_compra}\quad\tabla{p6_auc}\caption{Izquierda: fracción de clientes que
compró entre el 1 de octubre y el 8 de diciembre, por puntaje. Derecha: AUC con los valores crudos.}\end{table}
\begin{figure}[H]\centering\figura[0.7]{p6_tasa_compra}\caption{Tasa de recompra por puntaje.}\end{figure}

En total, el \cf{tasaCompra} de los clientes volvió a comprar. La dimensión que mejor separa por AUC es
\cf{mejorDim} (AUC de R: \cf{aucR}; F: \cf{aucF}; M: \cf{aucM}; suma de puntajes: \cf{aucSuma}). La tasa de recompra
crece con los tres puntajes, pero F tiene el mayor recorrido: del 32\,\% en el puntaje 1 al 86\,\% en el 5. Un
cliente que compra seguido es el que vuelve. R separa bien abajo (28\,\% contra 66\,\%), pero se satura: los
puntajes 4 y 5 compran igual, porque haber comprado hace 10 o hace 30 días da lo mismo en un horizonte de diez
semanas. En F, los puntajes 1 y 2 compran casi igual (32\,\% y 31\,\%) porque ambos son casi solo clientes con
$F=1$ repartidos al azar: es el desempate de 6.2 en acción. La suma de los tres puntajes es el mejor
clasificador, pero apenas mejora a F solo.

\subsection*{6.4 Campeones}

Los campeones (puntaje 4 o 5 en R, F y M) son \cf{nCampeones} clientes, el \cf{fracCampeones} de la base, y generaron
el \cf{shareCampeones} de la venta neta del 1 de octubre al 8 de diciembre de todos los clientes con
\texttt{CustomerID} (\cf{totalPost}), incluidos los nuevos.

\begin{table}[H]\centering\tabla{p6_grupos}\caption{Venta neta del 1 de octubre al 8 de diciembre por grupo.}\end{table}

\subsection*{6.5 Clientes nuevos}

\cf{nNuevos} clientes compraron por primera vez entre el 1 de octubre y el 8 de diciembre y generaron el
\cf{shareNuevos} de la venta neta del periodo. Ningún análisis RFM, ni ningún modelo construido cliente por cliente
con la historia hasta el 1 de octubre, puede pronosticar esa venta: esos clientes no existen en la tabla del modelo.
No tienen R, F ni M, y un modelo que suma pronósticos por cliente les asigna cero.

En el vocabulario de FPP2, la Parte 5 es un pronóstico \emph{top-down}: pronostica el total, y el total sí incluye a
los clientes nuevos, porque la llegada de clientes nuevos en el pasado ya está en la serie (por ejemplo, en la
subida de otoño). Un modelo cliente por cliente es \emph{bottom-up}: suma los pronósticos de cada cliente conocido y
pierde, por construcción, la venta de clientes que aún no conoce. Para el proyecto, donde compararemos ambos
enfoques, esto significa que el bottom-up tiene un sesgo hacia abajo de aproximadamente el \cf{shareNuevos} en este
periodo, a menos que se le agregue un componente aparte para clientes nuevos. A cambio, el bottom-up dice
\emph{quién} va a comprar, que es lo que el top-down no puede decir.

%======================================================================
\section*{Uso de herramientas de IA}
%======================================================================
Usamos un modelo de lenguaje (Claude) como apoyo para entender el enunciado, proponer la estructura del código y
redactar borradores de las explicaciones. Ejecutamos todo el código con los datos originales y revisamos cada
cifra e interpretación contra los resultados del notebook.

\begin{thebibliography}{9}
\bibitem{uci} D. Chen. \emph{Online Retail} [conjunto de datos]. UCI Machine Learning Repository, 2015.
\bibitem{chen} D. Chen, S. L. Sain y K. Guo. Data mining for the online retail industry: A case study of RFM
model-based customer segmentation using data mining. \emph{Journal of Database Marketing \& Customer Strategy
Management}, 19(3):197--208, 2012.
\bibitem{fpp2} R. J. Hyndman y G. Athanasopoulos. \emph{Forecasting: Principles and Practice}, 2a ed. OTexts, 2018.
\bibitem{tsay} R. S. Tsay. \emph{Analysis of Financial Time Series}, 2a ed. Wiley, 2005.
\end{thebibliography}

\end{document}
